% ==============================================================================
% SECTION 11: PRBS DATA GENERATION
% ==============================================================================
\section{Advanced Control Law}

\begin{frame}{Data generation}
    \framesubtitle{PRBS excitation in the 3D drone workspace}

    \begin{columns}[T,onlytextwidth]

        \begin{column}{0.48\textwidth}
            \begin{block}{Objective}
                Generate sufficiently rich input/output trajectories to
                identify a nonlinear dynamic model of the drone.
                \begin{itemize}
                    \item Excite the four actuator commands around hover
                    \item Capture the coupled 3D dynamics (12 variables)
                    \item Store synchronized inputs and measured outputs
                \end{itemize}
            \end{block}

            \begin{block}{Simulation setup}
                \begin{itemize}
                    \item Sampling time: $T_s = 0.05\,\mathrm{s}$
                    \item $N = 2000$ samples per trajectory
                    \item $3000$ independent experiments
                    \item Initial hover command: $U_{\mathrm{hover}} = 0.55$
                \end{itemize}
            \end{block}
        \end{column}

        \begin{column}{0.48\textwidth}
            \begin{block}{APRBS modes}
                We don't directly generate u inputs (unusable data).
                Instead, we use a flight controller : the inputs are now four independently generated piecewise-constant modes:
                \[
                    \delta = [\delta_T,\delta_r,\delta_p,\delta_y]
                \]

                Each level is sampled uniformly and held for a random duration
                between $1$ and $3$ seconds.
            \end{block}
        \end{column}
    \end{columns}
\end{frame}

\begin{frame}{Data generation}
    \framesubtitle{Flight controller}
    \begin{columns}[T,onlytextwidth]
        \begin{column}{0.50\textwidth}
            \begin{block}{From version 1 to version 2}
                The first version used PID controllers directly on the
                attitude angles. This solution was stable, but not fast
                enough for the drone dynamics (see appendix).

                Version 2 therefore uses a cascade controller: an outer
                angle loop generates an angular-rate reference, followed by
                a fast inner PID loop based on the gyroscope measurement.
            \end{block}

            \begin{block}{Controller parameters}
                \begin{itemize}
                    \item Outer angle gain: $K_a=2.5$
                    \item Inner rate PID:
                    $P=0.125,\ I=0.05,\ D=0.025,\ N=100$
                    \item Sampling time: $T_s=0.005\,\mathrm{s}$
                    \item Same structure for roll and pitch
                \end{itemize}
            \end{block}
        \end{column}

        \begin{column}{0.49\textwidth}
            \centering
            \resizebox{\linewidth}{!}{%
            \begin{tikzpicture}[
                blk/.style={draw=CSNavy,thick,fill=white,align=center,
                    inner sep=3pt,minimum height=8mm},
                sum/.style={circle,draw=CSNavy,thick,fill=white,
                    minimum size=5mm,inner sep=0pt},
                sig/.style={->,thick,CSNavy},
                rate/.style={->,thick,DarkGreen},
                angle/.style={->,thick,red!70!black},
                >={Stealth[length=2mm]},
                hop/.style={preaction={draw=white,line width=3pt}}]
                \node (ref) at (0,0) {$\phi_{\mathrm{ref}}$};
                \node[sum] (s1) at (1.2,0) {};
                \node[blk,minimum width=13mm] (ka) at (2.8,0)
                    {$K_a$\\[-2pt]\scriptsize $2.5$};
                \node[sum] (s2) at (4.4,0) {};
                \node[blk,minimum width=22mm] (pid) at (6.6,0)
                    {\scriptsize PID rate\\[-2pt]
                    \scriptsize $P/I/D=0.125/0.05/0.025$};
                \node[blk,minimum width=5mm,minimum height=25mm] (mux)
                    at (9.0,0) {\scriptsize Mux};
                \node[blk,minimum width=12mm,minimum height=25mm] (mix)
                    at (10.4,0) {$\bar M^{-1}$};
                \node[blk,minimum width=19mm,minimum height=25mm] (uav)
                    at (13.0,0) {\scriptsize UAV model};

                \draw[sig] (ref) -- (s1);
                \draw[sig] (s1) -- node[above,font=\scriptsize]{$e_\phi$} (ka);
                \draw[sig] (ka) -- node[above,font=\scriptsize]{$p_{\rm ref}$} (s2);
                \draw[sig] (s2) -- node[above,font=\scriptsize]{$e_p$} (pid);
                \draw[sig] (pid) -- node[above,font=\scriptsize]{$\delta_r$}
                    (mux);
                \draw[sig] (mux) -- (mix);
                \draw[sig] (mix) -- (uav);

                \draw[rate,hop] (13.9,-0.4) -- (13.9,-1.7) -- (4.4,-1.7)
                    -- (s2.south);
                \draw[angle,hop] (14.2,0.3) -- (14.2,-2.5) -- (1.2,-2.5)
                    -- (s1.south);
                \node[font=\scriptsize,DarkGreen] at (9,-1.7)
                    {fast inner rate loop};
                \node[font=\scriptsize,red!70!black] at (8,-2.5)
                    {outer angle loop};
            \end{tikzpicture}}
        \end{column}
    \end{columns}
\end{frame}
